\chapter{High-Performance Computing Implementation and Shared-State Analysis}
\label{chap:hpc_parallelization}

\section{Introduction and Parallelization Challenges in Phase-Field UELs}
\label{sec:hpc_intro}

Phase-field fracture modeling involves tightly coupled multi-field equations that must be solved across hundreds to thousands of non-linear equilibrium increments. To make fine-grid or adaptive simulations computationally tractable, high-performance computing (HPC) cluster architectures and multi-core parallel execution are indispensable.

In Abaqus/Standard, parallel solver execution is supported via two primary paradigms:
\begin{enumerate}
    \item \textbf{Shared-Memory Parallelism (OpenMP / Multi-Threading)}: Multiple execution threads share a single virtual address space on a shared-memory compute node (\texttt{mp\_mode=threads}).
    \item \textbf{Distributed-Memory Parallelism (MPI)}: Computational domains are partitioned across distinct processes with explicit message passing (\texttt{mp\_mode=mpi}).
\end{enumerate}

When custom user elements (UELs) interact with external databases or store shared state variables in Fortran \texttt{COMMON} blocks or \texttt{SAVE} arrays, multi-threaded execution introduces significant risk of race conditions, memory corruption, and non-deterministic numerical divergence. This chapter investigates the shared-state memory architecture, evaluates thread safety in user subroutines, and documents the fail-closed cluster execution framework.

\section{Abaqus UEXTERNALDB Common-Block Memory Architecture}
\label{sec:hpc_common_block_architecture}

\subsection{Inter-Subroutine Shared State Communication}
In the multi-layer UEL formulation (Section~\ref{sec:theory_uel_architecture}), user subroutine \texttt{UEXTERNALDB} acts as the persistent memory manager. At solver initialization (\texttt{LOP=0}), \texttt{UEXTERNALDB} allocates and populates global Fortran common blocks:
\begin{verbatim}
      COMMON /SV_ELEM_NODAL_PHASE_COM/ SV_ELEM_NODAL_PHASE(4, 100000)
      COMMON /SV_H_COMMITTED/          SV_H_COMMITTED(4, 100000)
\end{verbatim}
These arrays store the committed nodal phase values $d$ and integration-point history variables $\mathcal{H}$ across all elements.

During element stiffness assembly inside \texttt{UEL}, individual integration points access these common blocks to evaluate constitutive degradation and history updates:
\begin{equation}
\mathcal{H}_{\mathrm{old}} = \text{SV\_H\_COMMITTED}(k_{\mathrm{int}}, j_{\mathrm{elem}}).
\end{equation}

\subsection{Shared-State Safety in Multi-Threaded Execution}
When Abaqus/Standard executes under multi-threading (\texttt{threads=4}), element assembly loops are dynamically partitioned across worker threads. Because Fortran common blocks reside in global memory shared across all threads, concurrent read/write operations on shared element indices can produce severe race conditions unless strict thread isolation is maintained:
\begin{itemize}
    \item \textbf{Read-Only Shared Access}: If common blocks are populated once at \texttt{LOP=0} and treated as strictly read-only during increment iterations (as in the frozen history formulation), multiple threads can safely read simultaneously without locking.
    \item \textbf{Concurrent State Updates}: If worker threads write updated history $\mathcal{H}_{n+1}$ back into common arrays during element iterations, thread race conditions occur whenever neighboring elements or companion layers are processed asynchronously.
\end{itemize}

\section{Thread-Safety and Parallelization Boundary Audit}
\label{sec:hpc_parallel_audit}

\subsection{Controlled Rank and Thread Indexing}
To investigate thread identification inside Abaqus UEL callbacks, diagnostic probes were implemented utilizing Abaqus internal utilities \texttt{get\_thread\_id()} and \texttt{CALL GETRANK}:
\begin{itemize}
    \item \textbf{One-Rank / Four-Thread Evaluation (Stage D2C)}: A controlled single-notch transfer restart executed on 1 rank with 4 threads. The multi-threaded solution reproduced the serial reference solution exactly (relative difference $< 10^{-14}$), confirming that read-only common block access is thread-safe for the tested single-node configuration.
    \item \textbf{Dedicated Shared-State Benchmark (Task P3-T4)}: A dedicated multi-threaded shared-memory stress test was formulated to evaluate concurrent write safety. However, compilation of the test harness was blocked prior to linking due to compiler header interface mismatches.
\end{itemize}

\subsection{Formal Parallel Safety Claims and Boundaries}
In accordance with the thesis claim matrix:
\begin{enumerate}
    \item \textbf{Validated Claim}: One-rank multi-threaded execution with read-only common blocks is validated for the tested single-node configuration.
    \item \textbf{Withheld Claims}: General thread safety for concurrent read/write common blocks, distributed MPI safety, and hybrid MPI/OpenMP execution are \textbf{not supported} and remain strictly outside the validated scope of this thesis.
    \item \textbf{Production Execution Policy}: To guarantee 100\% numerical determinism and eliminate potential race conditions, all production adaptive validation runs are executed in single-core serial mode (1 CPU, 16\,GB RAM).
\end{enumerate}

\section{Cluster Execution Environment and Dual-Channel Notification System}
\label{sec:hpc_notification_architecture}

\subsection{HPC Environment Isolation and Toolchain Modules}
All cluster simulations were executed on the high-performance computing cluster at TU Bergakademie Freiberg using PBS Pro scheduler queues (\texttt{entry\_imfdfkmq} routed to \texttt{normal\_imfdfkmq}). The software environment enforces strict toolchain module isolation:
\begin{itemize}
    \item Fortran / C++ Compiler: Intel oneAPI 2024.2.0 (\texttt{module load intel/2024.2.0})
    \item GCC Compatibility Layer: GCC 11.4.0 (\texttt{module load gcc/11.4.0})
    \item Finite Element Solver: Abaqus 2023 (\texttt{module load abaqus/2023})
\end{itemize}

\subsection{Fail-Closed Dual-Channel Notification Design}
To guarantee real-time monitoring and complete provenance tracking across long-running cluster jobs, a fail-closed dual-channel notification protocol was designed and deployed (Fig.~\ref{fig:notification_architecture}):

\begin{figure}[htbp]
\centering
\small
\begin{verbatim}
+-------------------------------------------------------------------------+
|                  HPC Dual-Channel Notification Architecture             |
+-------------------------------------------------------------------------+
|  Channel 1: PBS Native System Mail                                      |
|    - Directives: #PBS -m abe (Abort, Begin, End)                        |
|    - Recipient: pr21vyci@mailserver.tu-freiberg.de                      |
|    - Fallback: mailx helper script (pbs_job_mail_notify.sh)             |
+-------------------------------------------------------------------------+
|  Channel 2: Telegram Bot Webhook API                                    |
|    - Credentials: ~/.config/hpc-notify/telegram.env (chmod 600)         |
|    - Script: telegram_notify.py (strict zero-leakage token handling)   |
|    - Lifecycle Hooks (pbs_notify.sh):                                   |
|        * SUBMITTED: Dispatched from login node upon qsub                |
|        * BEGIN:     Dispatched upon compute-node allocation start       |
|        * PASS/FAIL: Dispatched upon solver exit via trap handler        |
+-------------------------------------------------------------------------+
\end{verbatim}
\caption[Schematic of the dual-channel fail-closed HPC notification system.]{Schematic of the dual-channel fail-closed HPC notification system.}
\label{fig:notification_architecture}
\end{figure}

The notification workflow enforces:
\begin{itemize}
    \item \textbf{Security \& Zero Secret Leakage}: Bot tokens are stored exclusively in user-owned private files (\texttt{chmod 600}) and never exposed in command lines, PBS scripts, or Git logs.
    \item \textbf{Transport vs Receipt Separation}: Network transport success ($\text{return code} = 0$) confirms API dispatch only; human verification is logged independently.
\end{itemize}

\section{Chapter Summary}
\label{sec:hpc_summary}

The high-performance computing investigations establish:
\begin{enumerate}
    \item Read-only access to \texttt{UEXTERNALDB} common blocks is verified for one-rank multi-threaded execution, while concurrent write operations require careful synchronization.
    \item Production adaptive simulations are standardized on single-core serial execution to guarantee strict numerical reproducibility.
    \item The fail-closed dual-channel notification architecture ensures reliable tracking across all HPC job lifecycles.
\end{enumerate}
